% Options for packages loaded elsewhere
% Options for packages loaded elsewhere
\PassOptionsToPackage{unicode}{hyperref}
\PassOptionsToPackage{hyphens}{url}
\PassOptionsToPackage{dvipsnames,svgnames,x11names}{xcolor}
%
\documentclass[
]{article}
\usepackage{xcolor}
\usepackage[top=0.75in,left=0.75in,bottom=0.75in,right=0.75in]{geometry}
\usepackage{amsmath,amssymb}
\setcounter{secnumdepth}{5}
\usepackage{iftex}
\ifPDFTeX
  \usepackage[T1]{fontenc}
  \usepackage[utf8]{inputenc}
  \usepackage{textcomp} % provide euro and other symbols
\else % if luatex or xetex
  \usepackage{unicode-math} % this also loads fontspec
  \defaultfontfeatures{Scale=MatchLowercase}
  \defaultfontfeatures[\rmfamily]{Ligatures=TeX,Scale=1}
\fi
\usepackage{lmodern}
\ifPDFTeX\else
  % xetex/luatex font selection
\fi
% Use upquote if available, for straight quotes in verbatim environments
\IfFileExists{upquote.sty}{\usepackage{upquote}}{}
\IfFileExists{microtype.sty}{% use microtype if available
  \usepackage[]{microtype}
  \UseMicrotypeSet[protrusion]{basicmath} % disable protrusion for tt fonts
}{}
\makeatletter
\@ifundefined{KOMAClassName}{% if non-KOMA class
  \IfFileExists{parskip.sty}{%
    \usepackage{parskip}
  }{% else
    \setlength{\parindent}{0pt}
    \setlength{\parskip}{6pt plus 2pt minus 1pt}}
}{% if KOMA class
  \KOMAoptions{parskip=half}}
\makeatother
% Make \paragraph and \subparagraph free-standing
\makeatletter
\ifx\paragraph\undefined\else
  \let\oldparagraph\paragraph
  \renewcommand{\paragraph}{
    \@ifstar
      \xxxParagraphStar
      \xxxParagraphNoStar
  }
  \newcommand{\xxxParagraphStar}[1]{\oldparagraph*{#1}\mbox{}}
  \newcommand{\xxxParagraphNoStar}[1]{\oldparagraph{#1}\mbox{}}
\fi
\ifx\subparagraph\undefined\else
  \let\oldsubparagraph\subparagraph
  \renewcommand{\subparagraph}{
    \@ifstar
      \xxxSubParagraphStar
      \xxxSubParagraphNoStar
  }
  \newcommand{\xxxSubParagraphStar}[1]{\oldsubparagraph*{#1}\mbox{}}
  \newcommand{\xxxSubParagraphNoStar}[1]{\oldsubparagraph{#1}\mbox{}}
\fi
\makeatother

\usepackage{color}
\usepackage{fancyvrb}
\newcommand{\VerbBar}{|}
\newcommand{\VERB}{\Verb[commandchars=\\\{\}]}
\DefineVerbatimEnvironment{Highlighting}{Verbatim}{commandchars=\\\{\}}
% Add ',fontsize=\small' for more characters per line
\usepackage{framed}
\definecolor{shadecolor}{RGB}{241,243,245}
\newenvironment{Shaded}{\begin{snugshade}}{\end{snugshade}}
\newcommand{\AlertTok}[1]{\textcolor[rgb]{0.68,0.00,0.00}{#1}}
\newcommand{\AnnotationTok}[1]{\textcolor[rgb]{0.37,0.37,0.37}{#1}}
\newcommand{\AttributeTok}[1]{\textcolor[rgb]{0.40,0.45,0.13}{#1}}
\newcommand{\BaseNTok}[1]{\textcolor[rgb]{0.68,0.00,0.00}{#1}}
\newcommand{\BuiltInTok}[1]{\textcolor[rgb]{0.00,0.23,0.31}{#1}}
\newcommand{\CharTok}[1]{\textcolor[rgb]{0.13,0.47,0.30}{#1}}
\newcommand{\CommentTok}[1]{\textcolor[rgb]{0.37,0.37,0.37}{#1}}
\newcommand{\CommentVarTok}[1]{\textcolor[rgb]{0.37,0.37,0.37}{\textit{#1}}}
\newcommand{\ConstantTok}[1]{\textcolor[rgb]{0.56,0.35,0.01}{#1}}
\newcommand{\ControlFlowTok}[1]{\textcolor[rgb]{0.00,0.23,0.31}{\textbf{#1}}}
\newcommand{\DataTypeTok}[1]{\textcolor[rgb]{0.68,0.00,0.00}{#1}}
\newcommand{\DecValTok}[1]{\textcolor[rgb]{0.68,0.00,0.00}{#1}}
\newcommand{\DocumentationTok}[1]{\textcolor[rgb]{0.37,0.37,0.37}{\textit{#1}}}
\newcommand{\ErrorTok}[1]{\textcolor[rgb]{0.68,0.00,0.00}{#1}}
\newcommand{\ExtensionTok}[1]{\textcolor[rgb]{0.00,0.23,0.31}{#1}}
\newcommand{\FloatTok}[1]{\textcolor[rgb]{0.68,0.00,0.00}{#1}}
\newcommand{\FunctionTok}[1]{\textcolor[rgb]{0.28,0.35,0.67}{#1}}
\newcommand{\ImportTok}[1]{\textcolor[rgb]{0.00,0.46,0.62}{#1}}
\newcommand{\InformationTok}[1]{\textcolor[rgb]{0.37,0.37,0.37}{#1}}
\newcommand{\KeywordTok}[1]{\textcolor[rgb]{0.00,0.23,0.31}{\textbf{#1}}}
\newcommand{\NormalTok}[1]{\textcolor[rgb]{0.00,0.23,0.31}{#1}}
\newcommand{\OperatorTok}[1]{\textcolor[rgb]{0.37,0.37,0.37}{#1}}
\newcommand{\OtherTok}[1]{\textcolor[rgb]{0.00,0.23,0.31}{#1}}
\newcommand{\PreprocessorTok}[1]{\textcolor[rgb]{0.68,0.00,0.00}{#1}}
\newcommand{\RegionMarkerTok}[1]{\textcolor[rgb]{0.00,0.23,0.31}{#1}}
\newcommand{\SpecialCharTok}[1]{\textcolor[rgb]{0.37,0.37,0.37}{#1}}
\newcommand{\SpecialStringTok}[1]{\textcolor[rgb]{0.13,0.47,0.30}{#1}}
\newcommand{\StringTok}[1]{\textcolor[rgb]{0.13,0.47,0.30}{#1}}
\newcommand{\VariableTok}[1]{\textcolor[rgb]{0.07,0.07,0.07}{#1}}
\newcommand{\VerbatimStringTok}[1]{\textcolor[rgb]{0.13,0.47,0.30}{#1}}
\newcommand{\WarningTok}[1]{\textcolor[rgb]{0.37,0.37,0.37}{\textit{#1}}}

\usepackage{longtable,booktabs,array}
\newcounter{none} % for unnumbered tables
\usepackage{calc} % for calculating minipage widths
% Correct order of tables after \paragraph or \subparagraph
\usepackage{etoolbox}
\makeatletter
\patchcmd\longtable{\par}{\if@noskipsec\mbox{}\fi\par}{}{}
\makeatother
% Allow footnotes in longtable head/foot
\IfFileExists{footnotehyper.sty}{\usepackage{footnotehyper}}{\usepackage{footnote}}
\makesavenoteenv{longtable}
\usepackage{graphicx}
\makeatletter
\newsavebox\pandoc@box
\newcommand*\pandocbounded[1]{% scales image to fit in text height/width
  \sbox\pandoc@box{#1}%
  \Gscale@div\@tempa{\textheight}{\dimexpr\ht\pandoc@box+\dp\pandoc@box\relax}%
  \Gscale@div\@tempb{\linewidth}{\wd\pandoc@box}%
  \ifdim\@tempb\p@<\@tempa\p@\let\@tempa\@tempb\fi% select the smaller of both
  \ifdim\@tempa\p@<\p@\scalebox{\@tempa}{\usebox\pandoc@box}%
  \else\usebox{\pandoc@box}%
  \fi%
}
% Set default figure placement to htbp
\def\fps@figure{htbp}
\makeatother


% definitions for citeproc citations
\NewDocumentCommand\citeproctext{}{}
\NewDocumentCommand\citeproc{mm}{%
  \begingroup\def\citeproctext{#2}\cite{#1}\endgroup}
\makeatletter
 % allow citations to break across lines
 \let\@cite@ofmt\@firstofone
 % avoid brackets around text for \cite:
 \def\@biblabel#1{}
 \def\@cite#1#2{{#1\if@tempswa , #2\fi}}
\makeatother
\newlength{\cslhangindent}
\setlength{\cslhangindent}{1.5em}
\newlength{\csllabelwidth}
\setlength{\csllabelwidth}{3em}
\newenvironment{CSLReferences}[2] % #1 hanging-indent, #2 entry-spacing
 {\begin{list}{}{%
  \setlength{\itemindent}{0pt}
  \setlength{\leftmargin}{0pt}
  \setlength{\parsep}{0pt}
  % turn on hanging indent if param 1 is 1
  \ifodd #1
   \setlength{\leftmargin}{\cslhangindent}
   \setlength{\itemindent}{-1\cslhangindent}
  \fi
  % set entry spacing
  \setlength{\itemsep}{#2\baselineskip}}}
 {\end{list}}
\usepackage{calc}
\newcommand{\CSLBlock}[1]{\hfill\break\parbox[t]{\linewidth}{\strut\ignorespaces#1\strut}}
\newcommand{\CSLLeftMargin}[1]{\parbox[t]{\csllabelwidth}{\strut#1\strut}}
\newcommand{\CSLRightInline}[1]{\parbox[t]{\linewidth - \csllabelwidth}{\strut#1\strut}}
\newcommand{\CSLIndent}[1]{\hspace{\cslhangindent}#1}



\setlength{\emergencystretch}{3em} % prevent overfull lines

\providecommand{\tightlist}{%
  \setlength{\itemsep}{0pt}\setlength{\parskip}{0pt}}



 


\makeatletter
\@ifpackageloaded{caption}{}{\usepackage{caption}}
\AtBeginDocument{%
\ifdefined\contentsname
  \renewcommand*\contentsname{Table of contents}
\else
  \newcommand\contentsname{Table of contents}
\fi
\ifdefined\listfigurename
  \renewcommand*\listfigurename{List of Figures}
\else
  \newcommand\listfigurename{List of Figures}
\fi
\ifdefined\listtablename
  \renewcommand*\listtablename{List of Tables}
\else
  \newcommand\listtablename{List of Tables}
\fi
\ifdefined\figurename
  \renewcommand*\figurename{Figure}
\else
  \newcommand\figurename{Figure}
\fi
\ifdefined\tablename
  \renewcommand*\tablename{Table}
\else
  \newcommand\tablename{Table}
\fi
}
\@ifpackageloaded{float}{}{\usepackage{float}}
\floatstyle{ruled}
\@ifundefined{c@chapter}{\newfloat{codelisting}{h}{lop}}{\newfloat{codelisting}{h}{lop}[chapter]}
\floatname{codelisting}{Listing}
\newcommand*\listoflistings{\listof{codelisting}{List of Listings}}
\makeatother
\makeatletter
\makeatother
\makeatletter
\@ifpackageloaded{caption}{}{\usepackage{caption}}
\@ifpackageloaded{subcaption}{}{\usepackage{subcaption}}
\makeatother
\usepackage{bookmark}
\IfFileExists{xurl.sty}{\usepackage{xurl}}{} % add URL line breaks if available
\urlstyle{same}
% fallback for those not using the hyperref driver hyperxmp:
\makeatletter
\@ifundefined{xmpquote}{\newcommand{\xmpquote}[1]{#1}}{}
\makeatother
\hypersetup{
  pdftitle={Homework 3 -- Regression},
  colorlinks=true,
  linkcolor={blue},
  filecolor={Maroon},
  citecolor={Blue},
  urlcolor={Blue},
  pdfcreator={LaTeX via pandoc}}


\title{Homework 3 -- Regression}
\author{}
\date{}
\begin{document}
\maketitle


\normalsize

\normalsize

\section{Question 1}\label{question-1}

\textbf{Below, I simulate a rather quite annoying set of regression
data.}

\scriptsize

\begin{Shaded}
\begin{Highlighting}[numbers=left,,]
\FunctionTok{set.seed}\NormalTok{(}\DecValTok{0510}\NormalTok{)}
\NormalTok{n }\OtherTok{\textless{}{-}} \DecValTok{250}
\NormalTok{dat.q1 }\OtherTok{\textless{}{-}} \FunctionTok{tibble}\NormalTok{(}\AttributeTok{x1 =} \FunctionTok{rnorm}\NormalTok{(n, }\AttributeTok{mean =} \DecValTok{50}\NormalTok{, }\AttributeTok{sd =} \DecValTok{10}\NormalTok{),}
                 \AttributeTok{x2 =} \FunctionTok{factor}\NormalTok{(}\FunctionTok{sample}\NormalTok{(}\FunctionTok{c}\NormalTok{(}\StringTok{"A"}\NormalTok{, }\StringTok{"B"}\NormalTok{), n, }\AttributeTok{replace =} \ConstantTok{TRUE}\NormalTok{)),}
                 \AttributeTok{x3 =} \FunctionTok{factor}\NormalTok{(}\FunctionTok{sample}\NormalTok{(}\FunctionTok{c}\NormalTok{(}\StringTok{"Control"}\NormalTok{, }\StringTok{"Treat"}\NormalTok{), n, }\AttributeTok{replace =} \ConstantTok{TRUE}\NormalTok{)),}
                 \AttributeTok{x4 =} \FunctionTok{factor}\NormalTok{(}\FunctionTok{sample}\NormalTok{(}\FunctionTok{c}\NormalTok{(}\StringTok{"Low"}\NormalTok{, }\StringTok{"High"}\NormalTok{), n, }\AttributeTok{replace =} \ConstantTok{TRUE}\NormalTok{), }\AttributeTok{levels=}\FunctionTok{c}\NormalTok{(}\StringTok{"Low"}\NormalTok{, }\StringTok{"High"}\NormalTok{)),}
                 \AttributeTok{x5 =} \FunctionTok{rnorm}\NormalTok{(n, }\AttributeTok{mean =} \DecValTok{0}\NormalTok{, }\AttributeTok{sd =} \DecValTok{1}\NormalTok{),}
                 \AttributeTok{x6 =} \FunctionTok{factor}\NormalTok{(}\FunctionTok{sample}\NormalTok{(}\FunctionTok{c}\NormalTok{(}\StringTok{"Group1"}\NormalTok{, }\StringTok{"Group2"}\NormalTok{), n, }\AttributeTok{replace =} \ConstantTok{TRUE}\NormalTok{)),}
                 \AttributeTok{x7 =} \FunctionTok{rnorm}\NormalTok{(n, }\AttributeTok{mean =} \DecValTok{10}\NormalTok{, }\AttributeTok{sd =} \DecValTok{2}\NormalTok{),}
                 \AttributeTok{x8 =} \FunctionTok{rnorm}\NormalTok{(n, }\AttributeTok{mean =} \DecValTok{10}\NormalTok{, }\AttributeTok{sd =} \DecValTok{2}\NormalTok{)) }\SpecialCharTok{|\textgreater{}}
        \FunctionTok{mutate}\NormalTok{(}\AttributeTok{y =} \DecValTok{5} \SpecialCharTok{+}                          \CommentTok{\# Intercept}
                   \CommentTok{\# Quantiative x1}
                   \DecValTok{2} \SpecialCharTok{*}\NormalTok{ x1 }\SpecialCharTok{+}
                   \CommentTok{\# Categorical x2}
\NormalTok{                   (x2 }\SpecialCharTok{==} \StringTok{"B"}\NormalTok{) }\SpecialCharTok{*} \DecValTok{10} \SpecialCharTok{+}
                   \CommentTok{\# Categorical x Categorical (x3*x4)}
\NormalTok{                   (x3 }\SpecialCharTok{==} \StringTok{"Treat"}\NormalTok{) }\SpecialCharTok{*} \DecValTok{5} \SpecialCharTok{+}\NormalTok{ (x4 }\SpecialCharTok{==} \StringTok{"High"}\NormalTok{) }\SpecialCharTok{*} \DecValTok{2} \SpecialCharTok{+} 
\NormalTok{                   (x3 }\SpecialCharTok{==} \StringTok{"Treat"} \SpecialCharTok{\&}\NormalTok{ x4 }\SpecialCharTok{==} \StringTok{"High"}\NormalTok{) }\SpecialCharTok{*} \DecValTok{15} \SpecialCharTok{+}
                   \CommentTok{\# Quantitative x Categorical  (x5*x6)}
                   \DecValTok{3} \SpecialCharTok{*}\NormalTok{ x5 }\SpecialCharTok{+}\NormalTok{ (x6 }\SpecialCharTok{==} \StringTok{"Group2"}\NormalTok{) }\SpecialCharTok{*} \DecValTok{4} \SpecialCharTok{+} 
\NormalTok{                   (x5 }\SpecialCharTok{*}\NormalTok{ (x6 }\SpecialCharTok{==} \StringTok{"Group2"}\NormalTok{)) }\SpecialCharTok{*} \DecValTok{8} \SpecialCharTok{+} 
                   \CommentTok{\# Quantitative x Quantitative (x7*x8)}
                   \FloatTok{1.5} \SpecialCharTok{*}\NormalTok{ x7 }\SpecialCharTok{+} \FloatTok{0.5} \SpecialCharTok{*}\NormalTok{ x8 }\SpecialCharTok{+} 
                   \FloatTok{0.8} \SpecialCharTok{*}\NormalTok{ (x7 }\SpecialCharTok{*}\NormalTok{ x8) }\SpecialCharTok{+} 
                   \CommentTok{\# Error}
                   \FunctionTok{rnorm}\NormalTok{(n, }\AttributeTok{mean =} \DecValTok{0}\NormalTok{, }\AttributeTok{sd =} \DecValTok{15}\NormalTok{))}
\end{Highlighting}
\end{Shaded}

\normalsize

\subsection{Part a}\label{part-a}

\textbf{Fit the regression model.}

\begin{align*}
  Y_i =& \beta_0  +
       \beta_1 x_{1i} + \beta_2 I(x_{2i}=\text{B})+ \beta_3 I(x_{3i}=\text{Treat}) + \beta_4 I(x_{4i}=\text{High}) +\\
       & \beta_5 x_{5i} + \beta_6 I( x_{6i}=\text{Group2}) + \beta_7 x_{7i} + \beta_8 x_{8i}+\\
       & \beta_9 (I(x_{3i}=\text{Treat}) ~\times~ I(x_{4i}=\text{High})) + \tag{Categorical $\times$ Categorical}\\
        & \beta_{10} (x_{5i} ~\times~ I(x_{6i}=\text{Group2})) +\tag{Continuous $\times$ Categorical}\\
        & \beta_{11} (x_{7i} ~\times~ x_{8i}) +
         \epsilon_i \tag{Continuous $\times$ Continuous}
\end{align*}

\textbf{Solution: First, we fit the regression model using
\texttt{lm()}.}

\scriptsize

\begin{Shaded}
\begin{Highlighting}[numbers=left,,]
\NormalTok{mod.q1 }\OtherTok{\textless{}{-}} \FunctionTok{lm}\NormalTok{(y }\SpecialCharTok{\textasciitilde{}}\NormalTok{ x1 }\SpecialCharTok{+}\NormalTok{ x2 }\SpecialCharTok{+}\NormalTok{ x3}\SpecialCharTok{*}\NormalTok{x4 }\SpecialCharTok{+}\NormalTok{ x5}\SpecialCharTok{*}\NormalTok{x6 }\SpecialCharTok{+}\NormalTok{ x7}\SpecialCharTok{*}\NormalTok{x8, }\AttributeTok{data=}\NormalTok{dat.q1)}
\end{Highlighting}
\end{Shaded}

\normalsize

\subsection{Part b}\label{part-b}

\textbf{Conduct a residual and outlier analysis.}

\textbf{Residual Analysis}

\normalsize

\begin{Shaded}
\begin{Highlighting}[numbers=left,,]
\FunctionTok{plotResiduals}\NormalTok{(mod.q1)}
\end{Highlighting}
\end{Shaded}

\begin{figure}[H]

\caption{Residual diagnostics for the Question 1 model.}

{\centering \pandocbounded{\includegraphics[keepaspectratio]{HW3_files/figure-pdf/unnamed-chunk-4-1.pdf}}

}

\end{figure}%

\normalsize The residuals follow a Gaussian distribution and appear
balanced around 0. The residuals are independent via the simulated data.

\textbf{Outlier Analysis}

\normalsize

\begin{Shaded}
\begin{Highlighting}[numbers=left,,]
\FunctionTok{plotInfluence}\NormalTok{(mod.q1)}
\end{Highlighting}
\end{Shaded}

\begin{figure}[H]

\caption{Residual diagnostics for the Question 1 model.}

{\centering \pandocbounded{\includegraphics[keepaspectratio]{HW3_files/figure-pdf/unnamed-chunk-5-1.pdf}}

}

\end{figure}%

\normalsize There are very few outliers, and most only have a small
effect on the regression model.

\subsection{Part c}\label{part-c}

\textbf{Interpret the intercept, the coefficients on \(x_1\) and
\(x_2\).}

\(\hat{\beta}_0\): The expected value of \(\hat{Y} = 8.9653\) when the
quantitative variables \(x_1, x_5, x_7, x_8\) are 0 and when the
categorical variables \(x_2 = "A"\), \(x_3 = "Control"\),
\(x_4 = "Low"\), and \(x_6 = "Group1"\).

\(\hat{\beta}_1\): The expected change in \(\hat{Y} = 2.13546\) for a
one-unit increase in \(x_1\), all other variables constant.

\(\hat{\beta}_2\): The expected change in \(\hat{Y} = 10.02851\) for a
one-unit increase in \(x_2\), all other variables constant.

\subsection{Part d}\label{part-d}

\textbf{Untangle and interpret the \(x_3*x_4\) interaction using
algebra. Use Table~\ref{tbl-q1d-emmeans} and
Figure~\ref{fig-q1d-emmeans} for context.}

\scriptsize

\begin{Shaded}
\begin{Highlighting}[numbers=left,,]
\FunctionTok{library}\NormalTok{(emmeans)}
\NormalTok{emm.q1d }\OtherTok{\textless{}{-}} \FunctionTok{emmeans}\NormalTok{(mod.q1, }\SpecialCharTok{\textasciitilde{}}\NormalTok{x3}\SpecialCharTok{*}\NormalTok{x4) }\SpecialCharTok{|\textgreater{}} \CommentTok{\# Don\textquotesingle{}t average across levels: use |x2*x6}
            \FunctionTok{data.frame}\NormalTok{() }\SpecialCharTok{|\textgreater{}}
            \FunctionTok{rename}\NormalTok{(}\StringTok{"Mean"}  \OtherTok{=} \StringTok{"emmean"}\NormalTok{,}
                   \StringTok{"Lower"} \OtherTok{=} \StringTok{"lower.CL"}\NormalTok{,}
                   \StringTok{"Upper"} \OtherTok{=} \StringTok{"upper.CL"}\NormalTok{)}
\end{Highlighting}
\end{Shaded}

\normalsize

\normalsize

{

\begin{longtable}[]{@{}llrrrrr@{}}

\caption{\label{tbl-q1d-emmeans}Estimated marginal means for the model
in Question 1.}

\tabularnewline

\toprule\noalign{}
x3 & x4 & Mean & SE & df & Lower & Upper \\
\midrule\noalign{}
\endhead
\bottomrule\noalign{}
\endlastfoot
Control & Low & 212.9676 & 1.7871 & 238 & 209.4470 & 216.4882 \\
Treat & Low & 219.4792 & 1.9123 & 238 & 215.7120 & 223.2464 \\
Control & High & 216.4894 & 1.7916 & 238 & 212.9601 & 220.0188 \\
Treat & High & 238.1930 & 1.7631 & 238 & 234.7198 & 241.6662 \\

\end{longtable}

}

\normalsize

\scriptsize

\begin{Shaded}
\begin{Highlighting}[numbers=left,,]
\NormalTok{emm.q1d }\OtherTok{\textless{}{-}}\NormalTok{ emm.q1d }\SpecialCharTok{|\textgreater{}}
  \FunctionTok{mutate}\NormalTok{(}\AttributeTok{Observation =} \FunctionTok{paste}\NormalTok{(}\StringTok{"(x3="}\NormalTok{, x3, }\StringTok{", "}\NormalTok{, }\StringTok{"x4="}\NormalTok{, x4, }\StringTok{")"}\NormalTok{, }\AttributeTok{sep=}\StringTok{""}\NormalTok{))}
\FunctionTok{ggplot}\NormalTok{(emm.q1d) }\SpecialCharTok{+} 
  \FunctionTok{geom\_point}\NormalTok{(}\FunctionTok{aes}\NormalTok{(}\AttributeTok{x=}\NormalTok{Mean, }\AttributeTok{y=}\NormalTok{Observation))}\SpecialCharTok{+}
  \FunctionTok{geom\_errorbar}\NormalTok{(}\FunctionTok{aes}\NormalTok{(}\AttributeTok{y=}\NormalTok{Observation,}
                    \AttributeTok{xmin=}\NormalTok{Lower, }\AttributeTok{xmax=}\NormalTok{Upper),}
                \AttributeTok{width =} \FloatTok{0.25}\NormalTok{)}\SpecialCharTok{+}
  \CommentTok{\# facet\_wrap(\textasciitilde{}x2*x6)+}
  \FunctionTok{theme\_bw}\NormalTok{()}\SpecialCharTok{+}
  \FunctionTok{xlab}\NormalTok{(}\StringTok{"Estimated Mean of Y"}\NormalTok{)}
\end{Highlighting}
\end{Shaded}

\begin{figure}[H]

\caption{\label{fig-q1d-emmeans}Estimated marginal means for the model
in Question 1.}

\centering{

\pandocbounded{\includegraphics[keepaspectratio]{HW3_files/figure-pdf/fig-q1d-emmeans-1.pdf}}

}

\end{figure}%

\normalsize With all other predictors besides \(x_3\) and \(x_4\)
constant, the regression looks like:

\[E[Y_i] = C + \beta_3 I(x_{3i}=\text{Treat}) + \beta_4 I(x_{4i}=\text{High}) + \beta_9 (I(x_{3i}=\text{Treat}) ~\times~ I(x_{4i}=\text{High}))\]
When \(x_3 =\) ``Control'' and \(x_4 =\) ``Low'': \[ E[Y_i] = C \] When
\(x_3 =\) ``Treat'' and \(x_4 =\) ``Low'': \[ E[Y_i] = C + \beta_3\]
When \(x_3 =\) ``Control'' and \(x_4 =\) ``High'':
\[ E[Y_i] = C + \beta_4\] When \(x_3 =\) ``Treat'' and \(x_4 =\)
``High'': \[ E[Y_i] = C + \beta_3 + \beta_4 + \beta_9\]

All other predictors constant, when \(x_4 =\) ``Low'', \(Y_i\) is
expected to increase by \(\beta_3 = 6.51\) from Control to Treat.
\(x_4 =\) ``High'', \(Y_i\) is expected to increase by
\(\beta_3 + \beta_9 = 6.51 + 15.19 = 21.70\) from Control to Treat. The
effect of Treat increases as \(x_4\) moves from Low to High. When
\(x_3 =\) ``Control'', \(Y_i\) is expected to increase by
\(\beta_4 = 3.52\) from Low to High. When \(x_3 =\) ``Treat'', \(Y_i\)
is expected to increase by \(\beta_4 + \beta_9 = 3.52 + 15.19 = 18.71\)
from Low to High. The effect of High increases as \(x_3\) moves from
Control to Treat. These effects are confirmed through the table and the
figure.

\subsection{Part e}\label{part-e}

\textbf{Untangle and interpret the \(x_5*x_6\) interaction using
algebra. Use Table~\ref{tbl-q1e-emtrends} and
Figure~\ref{fig-q1e-emtrends} for context.}

\scriptsize

\begin{Shaded}
\begin{Highlighting}[numbers=left,,]
\NormalTok{emm.q1e }\OtherTok{\textless{}{-}} \FunctionTok{emtrends}\NormalTok{(mod.q1, }\AttributeTok{var=}\StringTok{"x5"}\NormalTok{, }\SpecialCharTok{\textasciitilde{}}\NormalTok{x5}\SpecialCharTok{|}\NormalTok{x6) }\SpecialCharTok{|\textgreater{}} 
            \FunctionTok{data.frame}\NormalTok{() }\SpecialCharTok{|\textgreater{}}
            \FunctionTok{rename}\NormalTok{(}\StringTok{"Slope"}  \OtherTok{=} \StringTok{"x5.trend"}\NormalTok{,}
                   \StringTok{"Lower"} \OtherTok{=} \StringTok{"lower.CL"}\NormalTok{,}
                   \StringTok{"Upper"} \OtherTok{=} \StringTok{"upper.CL"}\NormalTok{) }\SpecialCharTok{|\textgreater{}}
\NormalTok{            dplyr}\SpecialCharTok{::}\FunctionTok{select}\NormalTok{(x6, Slope, SE, df, Lower, Upper)}
\end{Highlighting}
\end{Shaded}

\normalsize

\normalsize

{

\begin{longtable}[]{@{}lrrrrr@{}}

\caption{\label{tbl-q1e-emtrends}Estimated marginal trends for the model
in Question 1.}

\tabularnewline

\toprule\noalign{}
x6 & Slope & SE & df & Lower & Upper \\
\midrule\noalign{}
\endhead
\bottomrule\noalign{}
\endlastfoot
Group1 & 2.4835 & 1.2942 & 238 & -0.0660 & 5.0330 \\
Group2 & 9.4309 & 1.3367 & 238 & 6.7977 & 12.0642 \\

\end{longtable}

}

\normalsize In Figure~\ref{fig-q1e-emtrends}, the slopes are the same
and just the intercepts vary depending on \(x_2\), \(x_3\), and \(x_4\).

\scriptsize

\begin{Shaded}
\begin{Highlighting}[numbers=left,,]
\NormalTok{emm.q1e }\OtherTok{\textless{}{-}} \FunctionTok{emmeans}\NormalTok{(mod.q1, }\SpecialCharTok{\textasciitilde{}}\NormalTok{x5}\SpecialCharTok{*}\NormalTok{x6,}
                    \AttributeTok{at =} \FunctionTok{list}\NormalTok{(}\AttributeTok{x5 =} \FunctionTok{seq}\NormalTok{(}\SpecialCharTok{{-}}\DecValTok{3}\NormalTok{, }\DecValTok{3}\NormalTok{, }\AttributeTok{length.out=}\DecValTok{500}\NormalTok{))) }\SpecialCharTok{|\textgreater{}} \CommentTok{\# Don\textquotesingle{}t average across levels: use |x2*x3*x4}
  \FunctionTok{data.frame}\NormalTok{() }\SpecialCharTok{|\textgreater{}}
            \FunctionTok{rename}\NormalTok{(}\StringTok{"Mean"}  \OtherTok{=} \StringTok{"emmean"}\NormalTok{,}
                   \StringTok{"Lower"} \OtherTok{=} \StringTok{"lower.CL"}\NormalTok{,}
                   \StringTok{"Upper"} \OtherTok{=} \StringTok{"upper.CL"}\NormalTok{)}

\FunctionTok{ggplot}\NormalTok{(emm.q1e) }\SpecialCharTok{+} 
  \FunctionTok{geom\_line}\NormalTok{(}\FunctionTok{aes}\NormalTok{(}\AttributeTok{x=}\NormalTok{x5, }\AttributeTok{y=}\NormalTok{Mean, }\AttributeTok{color=}\NormalTok{x6))}\SpecialCharTok{+}
  \FunctionTok{geom\_ribbon}\NormalTok{(}\FunctionTok{aes}\NormalTok{(}\AttributeTok{x=}\NormalTok{x5, }\AttributeTok{ymin=}\NormalTok{Lower, }\AttributeTok{ymax=}\NormalTok{Upper, }\AttributeTok{fill=}\NormalTok{x6), }\AttributeTok{alpha=}\FloatTok{0.25}\NormalTok{)}\SpecialCharTok{+}
  \CommentTok{\# facet\_wrap(\textasciitilde{}x2*x3*x4)+}
  \FunctionTok{theme\_bw}\NormalTok{()}
\end{Highlighting}
\end{Shaded}

\begin{figure}[H]

\caption{\label{fig-q1e-emtrends}Estimated marginal trends for the model
in Question 1.}

\centering{

\pandocbounded{\includegraphics[keepaspectratio]{HW3_files/figure-pdf/fig-q1e-emtrends-1.pdf}}

}

\end{figure}%

\normalsize

All other predictors constant, except \(x_5\) and \(x_6\), the
regression looks like:

\[ E[Y_i] = C + \beta_5 x_5 + \beta_6 I(x_{6i}=\text{Group2}) + \beta_{10} (x_5 \times I(x_{6i}=\text{Group2})) \]
Taking the partial derivative of \(E[Y_i]\) with respect to \(x_{5i}\)
gives us the slope:
\[\frac{\partial E[Y_i]}{\partial x_{5i}} = \beta_5 + \beta_{10}\, I(x_{6i}=\text{Group2})\]
When \(x_6 =\) ``Group1'', the slope of \(x_5 = \beta_5 = 2.48\). When
\(x_6 =\) ``Group2'', the slope of
\(x_5 = \beta_5 + \beta_{10} = 2.48 + 6.95 = 9.43\). These can be seen
in Table 2. Depending on the group, \(x_5\) has less (Group1) or more
(Group2) of an effect on \(Y\).

The difference between the two groups depends on \(x_5\). When \(x_6 =\)
``Group1'': \[E[Y_i] = C +\beta_5 x_{5i}\]. When \(x_6 =\) ``Group2'':
\[E[Y_i] = C + \beta_5 x_{5i} + \beta_6 + \beta_{10} x_{5i}\]. We can
see through Figure 2 that Group1 is higher than Group2 for low \(x_5\),
and vice versa for higher \(x_5\).

\subsection{Part f}\label{part-f}

Untangle and interpret the \(x_7*x_8\) interaction using algebra. Use
Table~\ref{tbl-q1f-emtrends-x7} and Table~\ref{tbl-q1f-emtrends-x8} for
context.

\scriptsize

\begin{Shaded}
\begin{Highlighting}[numbers=left,,]
\NormalTok{emm.q1f.x7 }\OtherTok{\textless{}{-}} \FunctionTok{emtrends}\NormalTok{(mod.q1, }\AttributeTok{var=}\StringTok{"x7"}\NormalTok{, }\SpecialCharTok{\textasciitilde{}}\NormalTok{x7}\SpecialCharTok{|}\NormalTok{x8, }
                    \AttributeTok{at =} \FunctionTok{list}\NormalTok{(}\AttributeTok{x8 =} \FunctionTok{mean}\NormalTok{(dat.q1}\SpecialCharTok{$}\NormalTok{x8) }\SpecialCharTok{+} \FunctionTok{c}\NormalTok{(}\SpecialCharTok{{-}}\DecValTok{1}\NormalTok{, }\DecValTok{0}\NormalTok{ ,}\DecValTok{1}\NormalTok{)}\SpecialCharTok{*}\FunctionTok{sd}\NormalTok{(dat.q1}\SpecialCharTok{$}\NormalTok{x8))) }\SpecialCharTok{|\textgreater{}} 
            \FunctionTok{data.frame}\NormalTok{() }\SpecialCharTok{|\textgreater{}}
            \FunctionTok{rename}\NormalTok{(}\StringTok{"Slope"}  \OtherTok{=} \StringTok{"x7.trend"}\NormalTok{,}
                   \StringTok{"Lower"} \OtherTok{=} \StringTok{"lower.CL"}\NormalTok{,}
                   \StringTok{"Upper"} \OtherTok{=} \StringTok{"upper.CL"}\NormalTok{) }\SpecialCharTok{|\textgreater{}}
\NormalTok{            dplyr}\SpecialCharTok{::}\FunctionTok{select}\NormalTok{(x8, Slope, SE, df, Lower, Upper)}
\end{Highlighting}
\end{Shaded}

\normalsize

\normalsize

{

\begin{longtable}[]{@{}rrrrrr@{}}

\caption{\label{tbl-q1f-emtrends-x7}Estimated marginal trends for the
model in Question 1.}

\tabularnewline

\toprule\noalign{}
x8 & Slope & SE & df & Lower & Upper \\
\midrule\noalign{}
\endhead
\bottomrule\noalign{}
\endlastfoot
8.1309 & 7.9230 & 0.6751 & 238 & 6.5932 & 9.2529 \\
10.0582 & 9.7129 & 0.5015 & 238 & 8.7250 & 10.7008 \\
11.9854 & 11.5027 & 0.7396 & 238 & 10.0458 & 12.9597 \\

\end{longtable}

}

\normalsize

\scriptsize

\begin{Shaded}
\begin{Highlighting}[numbers=left,,]
\NormalTok{emm.q1f.x8 }\OtherTok{\textless{}{-}} \FunctionTok{emtrends}\NormalTok{(mod.q1, }\AttributeTok{var=}\StringTok{"x8"}\NormalTok{, }\SpecialCharTok{\textasciitilde{}}\NormalTok{x8}\SpecialCharTok{|}\NormalTok{x7, }
                    \AttributeTok{at =} \FunctionTok{list}\NormalTok{(}\AttributeTok{x7 =} \FunctionTok{mean}\NormalTok{(dat.q1}\SpecialCharTok{$}\NormalTok{x7) }\SpecialCharTok{+} \FunctionTok{c}\NormalTok{(}\SpecialCharTok{{-}}\DecValTok{1}\NormalTok{, }\DecValTok{0}\NormalTok{ ,}\DecValTok{1}\NormalTok{)}\SpecialCharTok{*}\FunctionTok{sd}\NormalTok{(dat.q1}\SpecialCharTok{$}\NormalTok{x7))) }\SpecialCharTok{|\textgreater{}} 
            \FunctionTok{data.frame}\NormalTok{() }\SpecialCharTok{|\textgreater{}}
            \FunctionTok{rename}\NormalTok{(}\StringTok{"Slope"}  \OtherTok{=} \StringTok{"x8.trend"}\NormalTok{,}
                   \StringTok{"Lower"} \OtherTok{=} \StringTok{"lower.CL"}\NormalTok{,}
                   \StringTok{"Upper"} \OtherTok{=} \StringTok{"upper.CL"}\NormalTok{) }\SpecialCharTok{|\textgreater{}}
\NormalTok{            dplyr}\SpecialCharTok{::}\FunctionTok{select}\NormalTok{(x7, Slope, SE, df, Lower, Upper)}
\end{Highlighting}
\end{Shaded}

\normalsize

\normalsize

{

\begin{longtable}[]{@{}rrrrrr@{}}

\caption{\label{tbl-q1f-emtrends-x8}Estimated marginal trends for the
model in Question 1.}

\tabularnewline

\toprule\noalign{}
x7 & Slope & SE & df & Lower & Upper \\
\midrule\noalign{}
\endhead
\bottomrule\noalign{}
\endlastfoot
8.3824 & 6.8569 & 0.6069 & 238 & 5.6613 & 8.0525 \\
10.2208 & 8.5643 & 0.4744 & 238 & 7.6297 & 9.4989 \\
12.0592 & 10.2716 & 0.7325 & 238 & 8.8286 & 11.7146 \\

\end{longtable}

}

\normalsize

A Johnson Neyman analysis using the \texttt{interactions} package for
\texttt{R} (Long 2024), shows that the slope of \(x_7\) is statistically
discernible for all observed values of \(x_8\)
(Figure~\ref{fig-q1f-jni-x7}), and vice versa
(Figure~\ref{fig-q1f-jni-x8}).

\scriptsize

\begin{Shaded}
\begin{Highlighting}[numbers=left,,]
\FunctionTok{library}\NormalTok{(interactions)}
\FunctionTok{johnson\_neyman}\NormalTok{(mod.q1, }\AttributeTok{pred =}\NormalTok{ x7, }\AttributeTok{modx =}\NormalTok{ x8, }\AttributeTok{control.fdr =}\NormalTok{ T)}
\end{Highlighting}
\end{Shaded}

\begin{verbatim}
JOHNSON-NEYMAN INTERVAL

When x8 is OUTSIDE the interval [-13.08, 3.35], the slope of x7 is p < .05.

Note: The range of observed values of x8 is [4.69, 15.07]

Interval calculated using false discovery rate adjusted t = 1.97 
\end{verbatim}

\begin{figure}[H]

\caption{\label{fig-q1f-jni-x7}Johnson Neyman interval for the model in
Question 1.}

\centering{

\pandocbounded{\includegraphics[keepaspectratio]{HW3_files/figure-pdf/fig-q1f-jni-x7-1.pdf}}

}

\end{figure}%

\normalsize

\scriptsize

\begin{Shaded}
\begin{Highlighting}[numbers=left,,]
\FunctionTok{johnson\_neyman}\NormalTok{(mod.q1, }\AttributeTok{pred =}\NormalTok{ x8, }\AttributeTok{modx =}\NormalTok{ x7, }\AttributeTok{control.fdr =}\NormalTok{ T)}
\end{Highlighting}
\end{Shaded}

\begin{verbatim}
JOHNSON-NEYMAN INTERVAL

When x7 is OUTSIDE the interval [-9.96, 4.25], the slope of x8 is p < .05.

Note: The range of observed values of x7 is [4.89, 14.60]

Interval calculated using false discovery rate adjusted t = 1.97 
\end{verbatim}

\begin{figure}[H]

\caption{\label{fig-q1f-jni-x8}Johnson Neyman interval for the model in
Question 1.}

\centering{

\pandocbounded{\includegraphics[keepaspectratio]{HW3_files/figure-pdf/fig-q1f-jni-x8-1.pdf}}

}

\end{figure}%

\normalsize

All other predictors constant, except \(x_7\) and \(x_8\), the
regression looks like:

\[ E[Y_i] = C + \beta_7 x_{7i} + \beta_8 x_{8i} + \beta_{11} x_{7i} x_{8i} \]
Taking the partial derivative of \(E[Y_i]\) with respect to \(x_{7i}\)
and \(x_{8i}\) gives us the slopes:
\[\frac{\partial E[Y_i]}{\partial x_{7i}} = \beta_7 + \beta_{11} x_{8i}\]
\[\frac{\partial E[Y_i]}{\partial x_{8i}} = \beta_8 + \beta_{11} x_{7i}\]
In words, the effect of \(x_7\) on \(Y\) depends on \(x_8\) and vice
versa. Both slopes are linear functions of the other predictor. For each
one-unit increase in \(x_8\), the slope of \(x_7\) changes by
\(\beta_{11} = 0.93\), and vice versa, seen in Tables 3 and 4. Between
the first and last rows of Table 3, \(x_8\) increases by 3.85 and the
slope of \(x_7\) increases by 3.58, giving 3.58/3.85 = 0.93.

\(\beta_7 = 0.37\) represents the effect of \(x_7\) on \(Y\) when
\(x_8 = 0\), and \(\beta_8 = -0.93\) represents the effect of \(x_8\) on
\(Y\) when \(x_7 = 0\). However, since the range of both predictors
never reaches 0, these values aren't meaningful by themselves.

\subsection{Part g}\label{part-g}

Standardize the continuous data and refit the regression model. Redo
part f.~What do you notice about the coefficient estimates and why does
this make the untangling process easier?

\textbf{Solution:} First, we fit the regression model using
\texttt{lm()}.

\scriptsize

\begin{Shaded}
\begin{Highlighting}[numbers=left,,]
\CommentTok{\# Standardize Data}
\NormalTok{dat.z }\OtherTok{\textless{}{-}}\NormalTok{ dat.q1 }\SpecialCharTok{|\textgreater{}}
  \FunctionTok{mutate}\NormalTok{(}\FunctionTok{across}\NormalTok{(}\FunctionTok{c}\NormalTok{(x1, x5, x7, x8), }\SpecialCharTok{\textasciitilde{}} \FunctionTok{as.vector}\NormalTok{(}\FunctionTok{scale}\NormalTok{(.))))}
\CommentTok{\# Refit Model}
\NormalTok{mod.q1z }\OtherTok{\textless{}{-}} \FunctionTok{lm}\NormalTok{(y }\SpecialCharTok{\textasciitilde{}}\NormalTok{ x1 }\SpecialCharTok{+}\NormalTok{ x2 }\SpecialCharTok{+}\NormalTok{ x3}\SpecialCharTok{*}\NormalTok{x4 }\SpecialCharTok{+}\NormalTok{ x5}\SpecialCharTok{*}\NormalTok{x6 }\SpecialCharTok{+}\NormalTok{ x7}\SpecialCharTok{*}\NormalTok{x8, }\AttributeTok{data=}\NormalTok{dat.z)}
\end{Highlighting}
\end{Shaded}

\normalsize

\scriptsize

\begin{Shaded}
\begin{Highlighting}[numbers=left,,]
\NormalTok{emm.q1g.z7 }\OtherTok{\textless{}{-}} \FunctionTok{emtrends}\NormalTok{(mod.q1z, }\AttributeTok{var=}\StringTok{"x7"}\NormalTok{, }\SpecialCharTok{\textasciitilde{}}\NormalTok{x7}\SpecialCharTok{|}\NormalTok{x8,}
                    \AttributeTok{at =} \FunctionTok{list}\NormalTok{(}\AttributeTok{x8 =} \FunctionTok{c}\NormalTok{(}\SpecialCharTok{{-}}\DecValTok{1}\NormalTok{, }\DecValTok{0}\NormalTok{ ,}\DecValTok{1}\NormalTok{))) }\SpecialCharTok{|\textgreater{}}
          \FunctionTok{data.frame}\NormalTok{() }\SpecialCharTok{|\textgreater{}}
          \FunctionTok{rename}\NormalTok{(}\StringTok{"Slope"}  \OtherTok{=} \StringTok{"x7.trend"}\NormalTok{,}
                 \StringTok{"Lower"} \OtherTok{=} \StringTok{"lower.CL"}\NormalTok{,}
                 \StringTok{"Upper"} \OtherTok{=} \StringTok{"upper.CL"}\NormalTok{) }\SpecialCharTok{|\textgreater{}}
\NormalTok{          dplyr}\SpecialCharTok{::}\FunctionTok{select}\NormalTok{(x8, Slope, SE, df, Lower, Upper)}
\end{Highlighting}
\end{Shaded}

\normalsize

\normalsize

{

\begin{longtable}[]{@{}rrrrrr@{}}

\caption{\label{tbl-q1-emtrends-z7}Estimated marginal trends for the
model in Question 1 when quantitative variables are standardized.}

\tabularnewline

\toprule\noalign{}
x8 & Slope & SE & df & Lower & Upper \\
\midrule\noalign{}
\endhead
\bottomrule\noalign{}
\endlastfoot
-1 & 14.57 & 1.24 & 238 & 12.12 & 17.01 \\
0 & 17.86 & 0.92 & 238 & 16.04 & 19.67 \\
1 & 21.15 & 1.36 & 238 & 18.47 & 23.83 \\

\end{longtable}

}

\normalsize

\scriptsize

\begin{Shaded}
\begin{Highlighting}[numbers=left,,]
\NormalTok{emm.q1g.z8 }\OtherTok{\textless{}{-}} \FunctionTok{emtrends}\NormalTok{(mod.q1z, }\AttributeTok{var=}\StringTok{"x8"}\NormalTok{, }\SpecialCharTok{\textasciitilde{}}\NormalTok{x8}\SpecialCharTok{|}\NormalTok{x7,}
                      \AttributeTok{at =} \FunctionTok{list}\NormalTok{(}\AttributeTok{x7 =} \FunctionTok{c}\NormalTok{(}\SpecialCharTok{{-}}\DecValTok{1}\NormalTok{, }\DecValTok{0}\NormalTok{ ,}\DecValTok{1}\NormalTok{)))}\SpecialCharTok{|\textgreater{}}
              \FunctionTok{data.frame}\NormalTok{() }\SpecialCharTok{|\textgreater{}}
              \FunctionTok{rename}\NormalTok{(}\StringTok{"Slope"}  \OtherTok{=} \StringTok{"x8.trend"}\NormalTok{,}
                    \StringTok{"Lower"} \OtherTok{=} \StringTok{"lower.CL"}\NormalTok{,}
                    \StringTok{"Upper"} \OtherTok{=} \StringTok{"upper.CL"}\NormalTok{) }\SpecialCharTok{|\textgreater{}}
\NormalTok{              dplyr}\SpecialCharTok{::}\FunctionTok{select}\NormalTok{(x7, Slope, SE, df, Lower, Upper)}
\end{Highlighting}
\end{Shaded}

\normalsize

\normalsize

{

\begin{longtable}[]{@{}rrrrrr@{}}

\caption{\label{tbl-q1-emtrends-z8}Estimated marginal trends for the
model in Question 1 when quantitative variables are standardized.}

\tabularnewline

\toprule\noalign{}
x7 & Slope & SE & df & Lower & Upper \\
\midrule\noalign{}
\endhead
\bottomrule\noalign{}
\endlastfoot
-1 & 13.22 & 1.17 & 238 & 10.91 & 15.52 \\
0 & 16.51 & 0.91 & 238 & 14.70 & 18.31 \\
1 & 19.80 & 1.41 & 238 & 17.01 & 22.58 \\

\end{longtable}

}

\normalsize

AFter standardization, Taking the partial derivative of \(E[Y_i]\) with
respect to \(z_{7i}\) and \(z_{8i}\) gives us the slopes:
\[\frac{\partial E[Y_i]}{\partial z_{7i}} = \beta_7 + \beta_{11} z_{8i}\]
\[\frac{\partial E[Y_i]}{\partial z_{8i}} = \beta_8 + \beta_{11} z_{7i}\]

After standardizing both predictors, the values for
\(\beta_7, \beta_8, \beta_{11} = 17.86, 16.51, 3.29\), respectively.
Now, \(\beta_7\) represents the increase in \(Y\) after a one-SD
increase in \(x_7\) when \(x_8\) is at its average (\(z_8 = 0\)) and
\(\beta_8\) represents the increase in \(Y\) after a one-SD increase
\(x_8\) when \(x_7\) is at its average (\(z_7 = 0\)), due to the
standardization of the predictors around 0. This makes untangling easier
because we can see how the effect of each predictor behaves around the
average of the other predictor, which is more informative than knowing
the behavior of each predictor when the other predictor was 0.

\section{Question 2}\label{question-2}

Cui and Mao (2025) investigate how freezing and cryogenic temperatures
affect the strength, flexibility, stiffness, and internal breaking
behavior of 3D-printed CF/PEEK composites.

Carbon fiber reinforced PEEK (CF/PEEK) is a high-performance composite
used in 3D printing for aerospace, deep-sea, polar research. The
researchers aimed to assess whether and how extremely cold environments
impact this composite to understand how 3D-printed CF/PEEK performs in
freezing conditions. This work would be important to understand its
limits and develop preventative methods to protect against structural
failures.

The researchers created fifteen 3D printed dog-bone-shaped test
specimens and placed them in a chamber cooled with liquid nitrogen to
\(25^\circ\text{C}\), \(-25^\circ\text{C}\), \(-75^\circ\text{C}\),
\(-125^\circ\text{C}\), and \(-175^\circ\text{C}\). To test strength of
the specimen, a machine pulled both ends of the specimen until it
snapped.

In the table below, the researchers record the average tensile strength
for three specimens from each temperature level.

{\def\LTcaptype{none} % do not increment counter
\begin{longtable}[]{@{}ll@{}}
\toprule\noalign{}
Temperature & Mean Tensile Strength \\
\midrule\noalign{}
\endhead
\bottomrule\noalign{}
\endlastfoot
\(25^\circ\text{C}\) & 81.93 MPa \\
\(-25^\circ\text{C}\) & 76.63 MPa \\
\(-75^\circ\text{C}\) & 63.23 MPa \\
\(-125^\circ\text{C}\) & 32.13 MPa \\
\(-175^\circ\text{C}\) & 12.83 MPa \\
\end{longtable}
}

\subsection{Part a}\label{part-a-1}

Compute the estimates we derived using calculus in lecture 16:
\begin{align*} 
\widehat{\beta}_0 &= \bar{y} - \widehat{\beta}_1 \bar{x}_1 \\ 
\widehat{\beta}_1 &= \frac{\sum_{i=1}^n (x_{1i} - \bar{x}_1)(y_i-\bar{y})}{\sum_{i=1}^n (x_{1i}-\bar{x}_1)^2}.
\end{align*}

\normalsize

\begin{Shaded}
\begin{Highlighting}[numbers=left,,]
\CommentTok{\# data from table above}
\NormalTok{dat.q2 }\OtherTok{\textless{}{-}} \FunctionTok{tibble}\NormalTok{(}
  \AttributeTok{temp     =} \FunctionTok{c}\NormalTok{(}\DecValTok{25}\NormalTok{, }\SpecialCharTok{{-}}\DecValTok{25}\NormalTok{, }\SpecialCharTok{{-}}\DecValTok{75}\NormalTok{, }\SpecialCharTok{{-}}\DecValTok{125}\NormalTok{, }\SpecialCharTok{{-}}\DecValTok{175}\NormalTok{),          }\CommentTok{\# x (degrees C)}
  \AttributeTok{strength =} \FunctionTok{c}\NormalTok{(}\FloatTok{81.93}\NormalTok{, }\FloatTok{76.63}\NormalTok{, }\FloatTok{63.23}\NormalTok{, }\FloatTok{32.13}\NormalTok{, }\FloatTok{12.83}\NormalTok{)  }\CommentTok{\# y (MPa)}
\NormalTok{)}

\NormalTok{x }\OtherTok{\textless{}{-}}\NormalTok{ dat.q2}\SpecialCharTok{$}\NormalTok{temp}
\NormalTok{y }\OtherTok{\textless{}{-}}\NormalTok{ dat.q2}\SpecialCharTok{$}\NormalTok{strength}

\CommentTok{\# estimate of b1 (slope) using calculus}
\NormalTok{b1.hat }\OtherTok{\textless{}{-}} \FunctionTok{sum}\NormalTok{((x }\SpecialCharTok{{-}} \FunctionTok{mean}\NormalTok{(x)) }\SpecialCharTok{*}\NormalTok{ (y }\SpecialCharTok{{-}} \FunctionTok{mean}\NormalTok{(y))) }\SpecialCharTok{/} \FunctionTok{sum}\NormalTok{((x }\SpecialCharTok{{-}} \FunctionTok{mean}\NormalTok{(x))}\SpecialCharTok{\^{}}\DecValTok{2}\NormalTok{)}

\CommentTok{\# estimate of b0 (intercept) using calculus}
\NormalTok{b0.hat }\OtherTok{\textless{}{-}} \FunctionTok{mean}\NormalTok{(y) }\SpecialCharTok{{-}}\NormalTok{ b1.hat }\SpecialCharTok{*} \FunctionTok{mean}\NormalTok{(x)}

\NormalTok{b0.hat }\CommentTok{\# intercept}
\end{Highlighting}
\end{Shaded}

\begin{verbatim}
[1] 80.755
\end{verbatim}

\begin{Shaded}
\begin{Highlighting}[numbers=left,,]
\NormalTok{b1.hat }\CommentTok{\# slope}
\end{Highlighting}
\end{Shaded}

\begin{verbatim}
[1] 0.3654
\end{verbatim}

\normalsize

\subsection{Part b}\label{part-b-1}

Compute the estimates we derived using linear algebra in lecture 16:
\[\widehat{\boldsymbol{\beta}} = \left(\mathbf{X}'\mathbf{X}\right)^{-1} \mathbf{X}'\mathbf{Y}.\]
\textbf{Hint:} Run \texttt{?matrix()} in \texttt{R} for a quick reminder
about how to create a matrix.

\normalsize

\begin{Shaded}
\begin{Highlighting}[numbers=left,,]
\CommentTok{\# remake data from table above}
\NormalTok{dat.q2 }\OtherTok{\textless{}{-}} \FunctionTok{tibble}\NormalTok{(}
  \AttributeTok{temp     =} \FunctionTok{c}\NormalTok{(}\DecValTok{25}\NormalTok{, }\SpecialCharTok{{-}}\DecValTok{25}\NormalTok{, }\SpecialCharTok{{-}}\DecValTok{75}\NormalTok{, }\SpecialCharTok{{-}}\DecValTok{125}\NormalTok{, }\SpecialCharTok{{-}}\DecValTok{175}\NormalTok{),          }\CommentTok{\# x (degrees C)}
  \AttributeTok{strength =} \FunctionTok{c}\NormalTok{(}\FloatTok{81.93}\NormalTok{, }\FloatTok{76.63}\NormalTok{, }\FloatTok{63.23}\NormalTok{, }\FloatTok{32.13}\NormalTok{, }\FloatTok{12.83}\NormalTok{)  }\CommentTok{\# y (MPa)}
\NormalTok{)}

\NormalTok{x }\OtherTok{\textless{}{-}}\NormalTok{ dat.q2}\SpecialCharTok{$}\NormalTok{temp}
\NormalTok{y }\OtherTok{\textless{}{-}}\NormalTok{ dat.q2}\SpecialCharTok{$}\NormalTok{strength}

\CommentTok{\# make matrices X and Y}
\NormalTok{x.mat }\OtherTok{\textless{}{-}} \FunctionTok{cbind}\NormalTok{(}\DecValTok{1}\NormalTok{, x)}
\NormalTok{y.mat }\OtherTok{\textless{}{-}} \FunctionTok{matrix}\NormalTok{(}\AttributeTok{data =}\NormalTok{ y)}

\CommentTok{\# compute estimates using linear algebra}
\NormalTok{estimates }\OtherTok{\textless{}{-}} \FunctionTok{solve}\NormalTok{(}\FunctionTok{t}\NormalTok{(x.mat) }\SpecialCharTok{\%*\%}\NormalTok{ x.mat) }\SpecialCharTok{\%*\%} \FunctionTok{t}\NormalTok{(x.mat) }\SpecialCharTok{\%*\%}\NormalTok{ y.mat}
\NormalTok{b0.hat.b }\OtherTok{\textless{}{-}}\NormalTok{ estimates[}\DecValTok{1}\NormalTok{]}
\NormalTok{b1.hat.b }\OtherTok{\textless{}{-}}\NormalTok{ estimates[}\DecValTok{2}\NormalTok{]}

\NormalTok{b0.hat.b}
\end{Highlighting}
\end{Shaded}

\begin{verbatim}
[1] 80.755
\end{verbatim}

\begin{Shaded}
\begin{Highlighting}[numbers=left,,]
\NormalTok{b1.hat.b}
\end{Highlighting}
\end{Shaded}

\begin{verbatim}
[1] 0.3654
\end{verbatim}

\normalsize

\subsection{Part c}\label{part-c-1}

Compute the maximum likelihood estimates using \texttt{optim()} in
\texttt{R} as described in lecture 16:
\[\max_{\beta_0, \beta_1, \sigma} \mathcal{L}(\boldsymbol{\theta}=\{\beta_0, \beta_1, \sigma\}).\]

\normalsize

\begin{Shaded}
\begin{Highlighting}[numbers=left,,]
\CommentTok{\# Negative log{-}likelihood where theta = c(beta0, beta1, sigma)}
\NormalTok{neg.llike }\OtherTok{\textless{}{-}} \ControlFlowTok{function}\NormalTok{(theta, x, y) \{}
\NormalTok{  b0 }\OtherTok{\textless{}{-}}\NormalTok{    theta[}\DecValTok{1}\NormalTok{]}
\NormalTok{  b1 }\OtherTok{\textless{}{-}}\NormalTok{    theta[}\DecValTok{2}\NormalTok{]}
\NormalTok{  sigma }\OtherTok{\textless{}{-}}\NormalTok{ theta[}\DecValTok{3}\NormalTok{]}
\NormalTok{  mu }\OtherTok{\textless{}{-}}\NormalTok{ b0 }\SpecialCharTok{+}\NormalTok{ b1 }\SpecialCharTok{*}\NormalTok{ x                }\CommentTok{\# E(Y | x)}
  \SpecialCharTok{{-}}\FunctionTok{sum}\NormalTok{(}\FunctionTok{dnorm}\NormalTok{(y, }\AttributeTok{mean =}\NormalTok{ mu, }\AttributeTok{sd =}\NormalTok{ sigma, }\AttributeTok{log =} \ConstantTok{TRUE}\NormalTok{))}
\NormalTok{\}}

\NormalTok{mle }\OtherTok{\textless{}{-}} \FunctionTok{optim}\NormalTok{(}\AttributeTok{par =} \FunctionTok{c}\NormalTok{(}\FunctionTok{mean}\NormalTok{(y), }\DecValTok{0}\NormalTok{, }\FunctionTok{sd}\NormalTok{(y)),   }\CommentTok{\# starting values}
             \AttributeTok{fn  =}\NormalTok{ neg.llike,}
             \AttributeTok{x =}\NormalTok{ x, }\AttributeTok{y =}\NormalTok{ y,}
             \AttributeTok{control =} \FunctionTok{list}\NormalTok{(}\AttributeTok{maxit =} \DecValTok{10000}\NormalTok{, }\AttributeTok{reltol =} \FloatTok{1e{-}12}\NormalTok{))}

\NormalTok{mle }\CommentTok{\# ($par)}
\end{Highlighting}
\end{Shaded}

\begin{verbatim}
$par
[1] 80.754999  0.365400  6.490294

$value
[1] 16.44622

$counts
function gradient 
     190       NA 

$convergence
[1] 0

$message
NULL
\end{verbatim}

\normalsize

\subsection{Part d}\label{part-d-1}

Fit the model using \texttt{lm()} and plot the resulting model. Compare
to the results from parts (a) through (c).

\normalsize

\begin{Shaded}
\begin{Highlighting}[numbers=left,,]
\NormalTok{mod.q2 }\OtherTok{\textless{}{-}} \FunctionTok{lm}\NormalTok{(strength }\SpecialCharTok{\textasciitilde{}}\NormalTok{ temp, }\AttributeTok{data=}\NormalTok{dat.q2)}
\FunctionTok{summary}\NormalTok{(mod.q2)}
\end{Highlighting}
\end{Shaded}

\begin{verbatim}

Call:
lm(formula = strength ~ temp, data = dat.q2)

Residuals:
    1     2     3     4     5 
-7.96  5.01  9.88 -2.95 -3.98 

Coefficients:
            Estimate Std. Error t value Pr(>|t|)    
(Intercept) 80.75500    5.46239  14.784 0.000671 ***
temp         0.36540    0.05299   6.895 0.006250 ** 
---
Signif. codes:  0 '***' 0.001 '**' 0.01 '*' 0.05 '.' 0.1 ' ' 1

Residual standard error: 8.379 on 3 degrees of freedom
Multiple R-squared:  0.9406,    Adjusted R-squared:  0.9209 
F-statistic: 47.54 on 1 and 3 DF,  p-value: 0.00625
\end{verbatim}

\normalsize The results match the estimates found through parts (a) --
(c). \(\hat{\beta}_0 = 80.76\) (Intercept), and
\(\hat{\beta}_1 = 0.365\) (slope). In words, the strength at 0 degrees
Celsius is 80.76 MPa and the strength will decrease by 0.365 MPa with
every one-degree decrease in temperature.

\section{Question 3}\label{question-3}

{Bracht et al.} (2016) evaluate MFAP4 as a biomarker for severity of
fibrosis in hepatitic C patients.

\scriptsize

\begin{Shaded}
\begin{Highlighting}[numbers=left,,]
\NormalTok{dat.mfap4 }\OtherTok{\textless{}{-}} \FunctionTok{read\_csv}\NormalTok{(}\StringTok{"Bracht2016.csv"}\NormalTok{)}
\end{Highlighting}
\end{Shaded}

\normalsize

\subsection{Part a}\label{part-a-2}

Fit the model predicting \texttt{MFAP4} using \texttt{Fibrosis.Stage}.
Write out the resulting estimate model and interpret the coefficients.

\normalsize

\begin{Shaded}
\begin{Highlighting}[numbers=left,,]
\NormalTok{mod.q3a }\OtherTok{\textless{}{-}} \FunctionTok{lm}\NormalTok{(MFAP4 }\SpecialCharTok{\textasciitilde{}}\NormalTok{ Fibrosis.Stage, }\AttributeTok{data =}\NormalTok{ dat.mfap4)}
\FunctionTok{summary}\NormalTok{(mod.q3a)}
\end{Highlighting}
\end{Shaded}

\begin{verbatim}

Call:
lm(formula = MFAP4 ~ Fibrosis.Stage, data = dat.mfap4)

Residuals:
    Min      1Q  Median      3Q     Max 
-20.318  -5.518  -2.651   2.070  81.816 

Coefficients:
                      Estimate Std. Error t value Pr(>|t|)    
(Intercept)              9.230      1.102   8.378 4.73e-16 ***
Fibrosis.StageStage 1    1.301      1.372   0.948   0.3433    
Fibrosis.StageStage 2    3.154      1.444   2.184   0.0294 *  
Fibrosis.StageStage 3   11.742      1.724   6.812 2.59e-11 ***
Fibrosis.StageStage 4   15.388      1.724   8.928  < 2e-16 ***
---
Signif. codes:  0 '***' 0.001 '**' 0.01 '*' 0.05 '.' 0.1 ' ' 1

Residual standard error: 10.85 on 537 degrees of freedom
Multiple R-squared:  0.1965,    Adjusted R-squared:  0.1905 
F-statistic: 32.84 on 4 and 537 DF,  p-value: < 2.2e-16
\end{verbatim}

\normalsize

The resulting regression model:
\[ \hat{MFAP4_i} = 9.23 + 1.31 I(Fibrosis.Stage_i = \text{Stage 1}) + 3.15 I(Fibrosis.Stage_i = \text{Stage 2}) + 11.72 I(Fibrosis.Stage_i = \text{Stage 3}) + 15.39 I(Fibrosis.Stage_i = \text{Stage 4})\]
9.23 represents \(E[MFAP4_i]\) at Stage 0. 1.31 represents the increase
in \(E[MFAP4_i]\) from Stage 0 to Stage 1. 3.15 represents the increase
in \(E[MFAP4_i]\) from Stage 0 to Stage 2. 11.74 represents the increase
in \(E[MFAP4_i]\) from Stage 0 to Stage 3. 15.39 represents the increase
in \(E[MFAP4_i]\) from Stage 0 to Stage 4.

\subsection{Part b}\label{part-b-2}

Conduct a residual analysis on the model in part a to demonstrate how to
identify when you need a log transformation.

\normalsize

\begin{Shaded}
\begin{Highlighting}[numbers=left,,]
\FunctionTok{plotResiduals}\NormalTok{(mod.q3a)}
\end{Highlighting}
\end{Shaded}

\begin{figure}[H]

\caption{Residual diagnostics for the Question 3 model.}

{\centering \pandocbounded{\includegraphics[keepaspectratio]{HW3_files/figure-pdf/unnamed-chunk-19-1.pdf}}

}

\end{figure}%

\normalsize The residuals do not appear to be Gaussian distributed, the
values move away from the qq-plot as they increase. Many of the
residuals are clusted around 0 or below, with only positive outliers
present. This spread is conducive to a log transformation.

\subsection{Part c}\label{part-c-2}

Fit the model predicting \texttt{log(MFAP4)} using
\texttt{Fibrosis.Stage}. Use the natural log, and not the
\(\log_2(\cdot)\) recorded in the data file. Write out the resulting
estimate model and interpret the coefficients.

\normalsize

\begin{Shaded}
\begin{Highlighting}[numbers=left,,]
\NormalTok{mod.q3log }\OtherTok{\textless{}{-}} \FunctionTok{lm}\NormalTok{(}\FunctionTok{log}\NormalTok{(MFAP4) }\SpecialCharTok{\textasciitilde{}}\NormalTok{ Fibrosis.Stage, }\AttributeTok{data =}\NormalTok{ dat.mfap4)}
\FunctionTok{summary}\NormalTok{(mod.q3log)}
\end{Highlighting}
\end{Shaded}

\begin{verbatim}

Call:
lm(formula = log(MFAP4) ~ Fibrosis.Stage, data = dat.mfap4)

Residuals:
     Min       1Q   Median       3Q      Max 
-1.55668 -0.41861 -0.04282  0.35845  2.21571 

Coefficients:
                      Estimate Std. Error t value Pr(>|t|)    
(Intercept)            2.02572    0.06178  32.789  < 2e-16 ***
Fibrosis.StageStage 1  0.13336    0.07694   1.733 0.083637 .  
Fibrosis.StageStage 2  0.30399    0.08099   3.753 0.000193 ***
Fibrosis.StageStage 3  0.77408    0.09666   8.009 7.24e-15 ***
Fibrosis.StageStage 4  0.98957    0.09666  10.238  < 2e-16 ***
---
Signif. codes:  0 '***' 0.001 '**' 0.01 '*' 0.05 '.' 0.1 ' ' 1

Residual standard error: 0.6085 on 537 degrees of freedom
Multiple R-squared:  0.2312,    Adjusted R-squared:  0.2255 
F-statistic: 40.38 on 4 and 537 DF,  p-value: < 2.2e-16
\end{verbatim}

\normalsize The resulting regression model:
\[ \hat{MFAP4_i} = 2.03 + 0.13 I(Fibrosis.Stage_i = \text{Stage 1}) + 0.30 I(Fibrosis.Stage_i = \text{Stage 2}) + 0.77 I(Fibrosis.Stage_i = \text{Stage 3}) + 0.99 I(Fibrosis.Stage_i = \text{Stage 4})\]
2.03 represents \(E[ln(MFAP4_i)]\) at Stage 0. 0.13 represents the
increase in \(E[ln(MFAP4_i)]\) from Stage 0 to Stage 1. 0.30 represents
the increase in \(E[ln(MFAP4_i)]\) from Stage 0 to Stage 2. 0.77
represents the increase in \(E[ln(MFAP4_i)]\) from Stage 0 to Stage 3.
0.99 represents the increase in \(E[ln(MFAP4_i)]\) from Stage 0 to Stage
4.

\subsection{Part d}\label{part-d-2}

Conduct a residual analysis on the model in part c to demonstrate that
the conditions for using regression are met.

\normalsize

\begin{Shaded}
\begin{Highlighting}[numbers=left,,]
\FunctionTok{plotResiduals}\NormalTok{(mod.q3log)}
\end{Highlighting}
\end{Shaded}

\begin{figure}[H]

\caption{Residual diagnostics for the Question 3 model after log
transformation.}

{\centering \pandocbounded{\includegraphics[keepaspectratio]{HW3_files/figure-pdf/unnamed-chunk-21-1.pdf}}

}

\end{figure}%

\normalsize The residuals seem to be Gaussian distributed, and the
residuals are more evenly centered around 0, with very few outliers.
This is a better fit for a linear regression than the raw MFAP4 counts.

\subsection{Part e}\label{part-e-1}

Fit the model predicting \texttt{log(MFAP4)} using
\texttt{Fibrosis.Stage}, \texttt{Age}, and \texttt{Gender}. Use the
natural log, and not the \(\log_2(\cdot)\) recorded in the data file.
Write out the resulting estimate model and interpret the coefficients.

\normalsize

\begin{Shaded}
\begin{Highlighting}[numbers=left,,]
\NormalTok{mod.q3e }\OtherTok{\textless{}{-}} \FunctionTok{lm}\NormalTok{(}\FunctionTok{log}\NormalTok{(MFAP4) }\SpecialCharTok{\textasciitilde{}}\NormalTok{ Fibrosis.Stage }\SpecialCharTok{+}\NormalTok{ Age }\SpecialCharTok{+}\NormalTok{ Gender, }\AttributeTok{data =}\NormalTok{ dat.mfap4)}
\FunctionTok{summary}\NormalTok{(mod.q3e)}
\end{Highlighting}
\end{Shaded}

\begin{verbatim}

Call:
lm(formula = log(MFAP4) ~ Fibrosis.Stage + Age + Gender, data = dat.mfap4)

Residuals:
     Min       1Q   Median       3Q      Max 
-1.55799 -0.40010 -0.03682  0.35968  2.28903 

Coefficients:
                       Estimate Std. Error t value Pr(>|t|)    
(Intercept)            1.518667   0.115808  13.114  < 2e-16 ***
Fibrosis.StageStage 1  0.126858   0.073594   1.724  0.08533 .  
Fibrosis.StageStage 2  0.212539   0.079095   2.687  0.00743 ** 
Fibrosis.StageStage 3  0.598866   0.096575   6.201 1.12e-09 ***
Fibrosis.StageStage 4  0.828547   0.097208   8.523  < 2e-16 ***
Age                    0.013110   0.002123   6.175 1.31e-09 ***
Gendermale            -0.143405   0.050807  -2.823  0.00494 ** 
---
Signif. codes:  0 '***' 0.001 '**' 0.01 '*' 0.05 '.' 0.1 ' ' 1

Residual standard error: 0.5812 on 535 degrees of freedom
Multiple R-squared:  0.3012,    Adjusted R-squared:  0.2934 
F-statistic: 38.44 on 6 and 535 DF,  p-value: < 2.2e-16
\end{verbatim}

\normalsize The resulting regression model:
\[ \hat{MFAP4_i} = 1.52 +  I(Fibrosis.Stage_i = \text{Stage 1}) + 0.30 I(Fibrosis.Stage_i = \text{Stage 2}) + 0.77 I(Fibrosis.Stage_i = \text{Stage 3}) + 0.99 I(Fibrosis.Stage_i = \text{Stage 4})\]

2.03 represents \(E[ln(MFAP4_i)]\) at Stage 0. 0.13 represents the
increase in \(E[ln(MFAP4_i)]\) from Stage 0 to Stage 1. 0.21 represents
the increase in \(E[ln(MFAP4_i)]\) from Stage 0 to Stage 2. 0.60
represents the increase in \(E[ln(MFAP4_i)]\) from Stage 0 to Stage 3.
0.83 represents the increase in \(E[ln(MFAP4_i)]\) from Stage 0 to Stage
4. 0.01 represents the change in \(E[ln(MFAP4_i)]\) after a one-year
increase in age. -0.14 represents the change in \(E[ln(MFAP4)]\) from
female to male patients.

\section{Question 4}\label{question-4}

\subsection{Part a}\label{part-a-3}

Simulate data to match regression assumptions in a loop that runs
\(R=1000\) times. \begin{align*}
    x_i &\sim \text{Uniform}(a=0,b=10)\\
    \epsilon_i &\sim \text{Gaussian}(\mu=0, \sigma=1)\\
    y_i &= 2 + 1.5 x_i + \epsilon_i\\
    n &= 30
\end{align*} In each iteration, fit a regression line to the data and
store the coefficient estimates. Plot the resulting estimates. What do
you notice?

\normalsize

\begin{Shaded}
\begin{Highlighting}[numbers=left,,]
\FunctionTok{set.seed}\NormalTok{(}\DecValTok{1206}\NormalTok{)}
\NormalTok{R }\OtherTok{\textless{}{-}} \DecValTok{1000}
\NormalTok{n }\OtherTok{\textless{}{-}} \DecValTok{30}

\CommentTok{\# create a space for the estimates}
\NormalTok{estimates\_a }\OtherTok{\textless{}{-}} \FunctionTok{tibble}\NormalTok{(}\AttributeTok{b0 =} \FunctionTok{rep}\NormalTok{(}\ConstantTok{NA}\NormalTok{, R),}
                    \AttributeTok{b1 =} \FunctionTok{rep}\NormalTok{(}\ConstantTok{NA}\NormalTok{, R))}

\CommentTok{\# simulate the data}
\ControlFlowTok{for}\NormalTok{ (r }\ControlFlowTok{in} \DecValTok{1}\SpecialCharTok{:}\NormalTok{R) \{}
  \CommentTok{\# X \textasciitilde{} Uniform(a = 0, b = 10)}
\NormalTok{  dat.q4 }\OtherTok{\textless{}{-}} \FunctionTok{tibble}\NormalTok{(}\AttributeTok{x =} \FunctionTok{runif}\NormalTok{(n, }\AttributeTok{min =} \DecValTok{0}\NormalTok{, }\AttributeTok{max =} \DecValTok{10}\NormalTok{)) }\SpecialCharTok{|\textgreater{}}
    \FunctionTok{mutate}\NormalTok{(}\AttributeTok{y =} \DecValTok{2}             \CommentTok{\# Intercept}
           \SpecialCharTok{+} \FloatTok{1.5} \SpecialCharTok{*}\NormalTok{ x         }\CommentTok{\# Xi predictor}
           \SpecialCharTok{+} \FunctionTok{rnorm}\NormalTok{(n, }\DecValTok{0}\NormalTok{, }\DecValTok{1}\NormalTok{)) }\CommentTok{\# Gaussian distributed errors}
  
  \CommentTok{\# model}
\NormalTok{  mod.q4 }\OtherTok{\textless{}{-}} \FunctionTok{lm}\NormalTok{(y }\SpecialCharTok{\textasciitilde{}}\NormalTok{ x, }\AttributeTok{data =}\NormalTok{ dat.q4)}
  
  \CommentTok{\# store estimates}
\NormalTok{  estimates\_a}\SpecialCharTok{$}\NormalTok{b0[r] }\OtherTok{\textless{}{-}} \FunctionTok{coef}\NormalTok{(mod.q4)[}\DecValTok{1}\NormalTok{] }\CommentTok{\# b0}
\NormalTok{  estimates\_a}\SpecialCharTok{$}\NormalTok{b1[r] }\OtherTok{\textless{}{-}} \FunctionTok{coef}\NormalTok{(mod.q4)[}\DecValTok{2}\NormalTok{] }\CommentTok{\# b1}
\NormalTok{\}}

\CommentTok{\# plot the estimates}
\NormalTok{p.b0.a }\OtherTok{\textless{}{-}} \FunctionTok{ggplot}\NormalTok{(estimates\_a, }\FunctionTok{aes}\NormalTok{(}\AttributeTok{x =}\NormalTok{ b0)) }\SpecialCharTok{+}
  \FunctionTok{geom\_histogram}\NormalTok{(}\FunctionTok{aes}\NormalTok{(}\AttributeTok{y =} \FunctionTok{after\_stat}\NormalTok{(density)), }\AttributeTok{bins =} \DecValTok{30}\NormalTok{,}
                 \AttributeTok{fill =} \StringTok{"lightgrey"}\NormalTok{, }\AttributeTok{color =} \StringTok{"black"}\NormalTok{) }\SpecialCharTok{+}
  \FunctionTok{geom\_density}\NormalTok{(}\AttributeTok{color =} \StringTok{"goldenrod"}\NormalTok{, }\AttributeTok{linewidth =} \DecValTok{1}\NormalTok{) }\SpecialCharTok{+}
  \FunctionTok{geom\_vline}\NormalTok{(}\AttributeTok{xintercept =} \DecValTok{2}\NormalTok{, }\AttributeTok{color =} \StringTok{"skyblue3"}\NormalTok{,}
             \AttributeTok{linetype =} \StringTok{"dashed"}\NormalTok{, }\AttributeTok{linewidth =} \DecValTok{1}\NormalTok{) }\SpecialCharTok{+}       \CommentTok{\# true beta\_0}
  \FunctionTok{labs}\NormalTok{(}\AttributeTok{x =} \FunctionTok{expression}\NormalTok{(}\FunctionTok{hat}\NormalTok{(beta)[}\DecValTok{0}\NormalTok{]), }\AttributeTok{y =} \StringTok{"Density"}\NormalTok{) }\SpecialCharTok{+}
  \FunctionTok{theme\_bw}\NormalTok{()}

\NormalTok{p.b1.a }\OtherTok{\textless{}{-}} \FunctionTok{ggplot}\NormalTok{(estimates\_a, }\FunctionTok{aes}\NormalTok{(}\AttributeTok{x =}\NormalTok{ b1)) }\SpecialCharTok{+}
  \FunctionTok{geom\_histogram}\NormalTok{(}\FunctionTok{aes}\NormalTok{(}\AttributeTok{y =} \FunctionTok{after\_stat}\NormalTok{(density)), }\AttributeTok{bins =} \DecValTok{30}\NormalTok{,}
                 \AttributeTok{fill =} \StringTok{"lightgrey"}\NormalTok{, }\AttributeTok{color =} \StringTok{"black"}\NormalTok{) }\SpecialCharTok{+}
  \FunctionTok{geom\_density}\NormalTok{(}\AttributeTok{color =} \StringTok{"goldenrod"}\NormalTok{, }\AttributeTok{linewidth =} \DecValTok{1}\NormalTok{) }\SpecialCharTok{+}
  \FunctionTok{geom\_vline}\NormalTok{(}\AttributeTok{xintercept =} \FloatTok{1.5}\NormalTok{, }\AttributeTok{color =} \StringTok{"skyblue3"}\NormalTok{,}
             \AttributeTok{linetype =} \StringTok{"dashed"}\NormalTok{, }\AttributeTok{linewidth =} \DecValTok{1}\NormalTok{) }\SpecialCharTok{+}       \CommentTok{\# true beta\_1}
  \FunctionTok{labs}\NormalTok{(}\AttributeTok{x =} \FunctionTok{expression}\NormalTok{(}\FunctionTok{hat}\NormalTok{(beta)[}\DecValTok{1}\NormalTok{]), }\AttributeTok{y =} \StringTok{"Density"}\NormalTok{) }\SpecialCharTok{+}
  \FunctionTok{theme\_bw}\NormalTok{()}

\NormalTok{p.b0.a }\SpecialCharTok{|}\NormalTok{ p.b1.a}
\end{Highlighting}
\end{Shaded}

\pandocbounded{\includegraphics[keepaspectratio]{HW3_files/figure-pdf/unnamed-chunk-23-1.pdf}}

\normalsize The values for \hat{\beta}\_0 and \hat{\beta}\_1 are
approximately normally distributed around their true values of 2 and
1.5, respectively. The spread of values for \hat{\beta}\_0 is larger
than the spread of values around \hat{\beta}\_1.

\subsection{Part b}\label{part-b-3}

Simulate data to match regression assumptions in a loop that runs
\(R=1000\) times. \begin{align*}
    x_i &\sim \text{Uniform}(a=0,b=10)\\
    \epsilon_i &\sim \text{Shifted Exponential}(\lambda=1, \mu=0)\\
    y_i &= 2 + 1.5 x_i + \epsilon_i\\
    n &= 30
\end{align*} In each iteration, fit a regression line to the data and
store the coefficient estimates. Plot the resulting estimates. What do
you notice?

\textbf{Hint:} Use the \texttt{tolerance::r2exp()} function with
\texttt{rate=1} and \texttt{shift=-1}.

\normalsize

\begin{Shaded}
\begin{Highlighting}[numbers=left,,]
\FunctionTok{set.seed}\NormalTok{(}\DecValTok{1206}\NormalTok{)}
\NormalTok{R }\OtherTok{\textless{}{-}} \DecValTok{1000}
\NormalTok{n }\OtherTok{\textless{}{-}} \DecValTok{30}

\CommentTok{\# create a space for the estimates}
\NormalTok{estimates\_b }\OtherTok{\textless{}{-}} \FunctionTok{tibble}\NormalTok{(}\AttributeTok{b0 =} \FunctionTok{rep}\NormalTok{(}\ConstantTok{NA}\NormalTok{, R),}
                    \AttributeTok{b1 =} \FunctionTok{rep}\NormalTok{(}\ConstantTok{NA}\NormalTok{, R))}

\CommentTok{\# simulate the data}
\ControlFlowTok{for}\NormalTok{ (r }\ControlFlowTok{in} \DecValTok{1}\SpecialCharTok{:}\NormalTok{R) \{}
  \CommentTok{\# X \textasciitilde{} Uniform(a = 0, b = 10)}
\NormalTok{  dat.q4 }\OtherTok{\textless{}{-}} \FunctionTok{tibble}\NormalTok{(}\AttributeTok{x =} \FunctionTok{runif}\NormalTok{(n, }\AttributeTok{min =} \DecValTok{0}\NormalTok{, }\AttributeTok{max =} \DecValTok{10}\NormalTok{)) }\SpecialCharTok{|\textgreater{}}
    \FunctionTok{mutate}\NormalTok{(}\AttributeTok{y =} \DecValTok{2}             \CommentTok{\# Intercept}
           \SpecialCharTok{+} \FloatTok{1.5} \SpecialCharTok{*}\NormalTok{ x         }\CommentTok{\# Xi predictor}
           \CommentTok{\# Shifted exponentially distributed errors}
           \SpecialCharTok{+} \FunctionTok{r2exp}\NormalTok{(n, }\AttributeTok{rate =} \DecValTok{1}\NormalTok{, }\AttributeTok{shift =} \SpecialCharTok{{-}}\DecValTok{1}\NormalTok{)) }
  
  \CommentTok{\# model}
\NormalTok{  mod.q4 }\OtherTok{\textless{}{-}} \FunctionTok{lm}\NormalTok{(y }\SpecialCharTok{\textasciitilde{}}\NormalTok{ x, }\AttributeTok{data =}\NormalTok{ dat.q4)}
  
  \CommentTok{\# store estimates}
\NormalTok{  estimates\_b}\SpecialCharTok{$}\NormalTok{b0[r] }\OtherTok{\textless{}{-}} \FunctionTok{coef}\NormalTok{(mod.q4)[}\DecValTok{1}\NormalTok{] }\CommentTok{\# b0}
\NormalTok{  estimates\_b}\SpecialCharTok{$}\NormalTok{b1[r] }\OtherTok{\textless{}{-}} \FunctionTok{coef}\NormalTok{(mod.q4)[}\DecValTok{2}\NormalTok{] }\CommentTok{\# b1}
\NormalTok{\}}

\CommentTok{\# plot the estimates}
\NormalTok{p.b0.b }\OtherTok{\textless{}{-}} \FunctionTok{ggplot}\NormalTok{(estimates\_b, }\FunctionTok{aes}\NormalTok{(}\AttributeTok{x =}\NormalTok{ b0)) }\SpecialCharTok{+}
  \FunctionTok{geom\_histogram}\NormalTok{(}\FunctionTok{aes}\NormalTok{(}\AttributeTok{y =} \FunctionTok{after\_stat}\NormalTok{(density)), }\AttributeTok{bins =} \DecValTok{30}\NormalTok{,}
                 \AttributeTok{fill =} \StringTok{"lightgrey"}\NormalTok{, }\AttributeTok{color =} \StringTok{"black"}\NormalTok{) }\SpecialCharTok{+}
  \FunctionTok{geom\_density}\NormalTok{(}\AttributeTok{color =} \StringTok{"goldenrod"}\NormalTok{, }\AttributeTok{linewidth =} \DecValTok{1}\NormalTok{) }\SpecialCharTok{+}
  \FunctionTok{geom\_vline}\NormalTok{(}\AttributeTok{xintercept =} \DecValTok{2}\NormalTok{, }\AttributeTok{color =} \StringTok{"skyblue3"}\NormalTok{,}
             \AttributeTok{linetype =} \StringTok{"dashed"}\NormalTok{, }\AttributeTok{linewidth =} \DecValTok{1}\NormalTok{) }\SpecialCharTok{+}       \CommentTok{\# true beta\_0}
  \FunctionTok{labs}\NormalTok{(}\AttributeTok{x =} \FunctionTok{expression}\NormalTok{(}\FunctionTok{hat}\NormalTok{(beta)[}\DecValTok{0}\NormalTok{]), }\AttributeTok{y =} \StringTok{"Density"}\NormalTok{) }\SpecialCharTok{+}
  \FunctionTok{theme\_bw}\NormalTok{()}

\NormalTok{p.b1.b }\OtherTok{\textless{}{-}} \FunctionTok{ggplot}\NormalTok{(estimates\_b, }\FunctionTok{aes}\NormalTok{(}\AttributeTok{x =}\NormalTok{ b1)) }\SpecialCharTok{+}
  \FunctionTok{geom\_histogram}\NormalTok{(}\FunctionTok{aes}\NormalTok{(}\AttributeTok{y =} \FunctionTok{after\_stat}\NormalTok{(density)), }\AttributeTok{bins =} \DecValTok{30}\NormalTok{,}
                 \AttributeTok{fill =} \StringTok{"lightgrey"}\NormalTok{, }\AttributeTok{color =} \StringTok{"black"}\NormalTok{) }\SpecialCharTok{+}
  \FunctionTok{geom\_density}\NormalTok{(}\AttributeTok{color =} \StringTok{"goldenrod"}\NormalTok{, }\AttributeTok{linewidth =} \DecValTok{1}\NormalTok{) }\SpecialCharTok{+}
  \FunctionTok{geom\_vline}\NormalTok{(}\AttributeTok{xintercept =} \FloatTok{1.5}\NormalTok{, }\AttributeTok{color =} \StringTok{"skyblue3"}\NormalTok{,}
             \AttributeTok{linetype =} \StringTok{"dashed"}\NormalTok{, }\AttributeTok{linewidth =} \DecValTok{1}\NormalTok{) }\SpecialCharTok{+}       \CommentTok{\# true beta\_1}
  \FunctionTok{labs}\NormalTok{(}\AttributeTok{x =} \FunctionTok{expression}\NormalTok{(}\FunctionTok{hat}\NormalTok{(beta)[}\DecValTok{1}\NormalTok{]), }\AttributeTok{y =} \StringTok{"Density"}\NormalTok{) }\SpecialCharTok{+}
  \FunctionTok{theme\_bw}\NormalTok{()}

\NormalTok{p.b0.b }\SpecialCharTok{|}\NormalTok{ p.b1.b}
\end{Highlighting}
\end{Shaded}

\pandocbounded{\includegraphics[keepaspectratio]{HW3_files/figure-pdf/unnamed-chunk-24-1.pdf}}

\normalsize The values for \hat{\beta}\_0 and \hat{\beta}\_1 are
approximately normally distributed around their true values of 2 and
1.5, respectively. Again, the spread of values for \hat{\beta}\_0 is
larger than the spread of values around \hat{\beta}\_1. The plots look
very similar to the plots from part a.

\subsection{Part c}\label{part-c-3}

Simulate data to match regression assumptions in a loop that runs
\(R=1000\) times. \begin{align*}
    x_i &\sim \text{Uniform}(a=0,b=10)\\
    \epsilon_i &\sim \text{Cauchy}(m=0, s=1)\\
    y_i &= 2 + 1.5 x_i + \epsilon_i\\
    n &= 30
\end{align*} In each iteration, fit a regression line to the data and
store the coefficient estimates. Plot the resulting estimates. What do
you notice?

\normalsize

\begin{Shaded}
\begin{Highlighting}[numbers=left,,]
\FunctionTok{set.seed}\NormalTok{(}\DecValTok{1206}\NormalTok{)}
\NormalTok{R }\OtherTok{\textless{}{-}} \DecValTok{1000}
\NormalTok{n }\OtherTok{\textless{}{-}} \DecValTok{30}

\CommentTok{\# create a space for the estimates}
\NormalTok{estimates\_c }\OtherTok{\textless{}{-}} \FunctionTok{tibble}\NormalTok{(}\AttributeTok{b0 =} \FunctionTok{rep}\NormalTok{(}\ConstantTok{NA}\NormalTok{, R),}
                    \AttributeTok{b1 =} \FunctionTok{rep}\NormalTok{(}\ConstantTok{NA}\NormalTok{, R))}

\CommentTok{\# simulate the data}
\ControlFlowTok{for}\NormalTok{ (r }\ControlFlowTok{in} \DecValTok{1}\SpecialCharTok{:}\NormalTok{R) \{}
  \CommentTok{\# X \textasciitilde{} Uniform(a = 0, b = 10)}
\NormalTok{  dat.q4 }\OtherTok{\textless{}{-}} \FunctionTok{tibble}\NormalTok{(}\AttributeTok{x =} \FunctionTok{runif}\NormalTok{(n, }\AttributeTok{min =} \DecValTok{0}\NormalTok{, }\AttributeTok{max =} \DecValTok{10}\NormalTok{)) }\SpecialCharTok{|\textgreater{}}
    \FunctionTok{mutate}\NormalTok{(}\AttributeTok{y =} \DecValTok{2}             \CommentTok{\# Intercept}
           \SpecialCharTok{+} \FloatTok{1.5} \SpecialCharTok{*}\NormalTok{ x         }\CommentTok{\# Xi predictor}
           \CommentTok{\# Cauchy distributed errors}
           \SpecialCharTok{+} \FunctionTok{rcauchy}\NormalTok{(n, }\AttributeTok{location =} \DecValTok{0}\NormalTok{, }\AttributeTok{scale =} \DecValTok{1}\NormalTok{)) }
  
  \CommentTok{\# model}
\NormalTok{  mod.q4 }\OtherTok{\textless{}{-}} \FunctionTok{lm}\NormalTok{(y }\SpecialCharTok{\textasciitilde{}}\NormalTok{ x, }\AttributeTok{data =}\NormalTok{ dat.q4)}
  
  \CommentTok{\# store estimates}
\NormalTok{  estimates\_c}\SpecialCharTok{$}\NormalTok{b0[r] }\OtherTok{\textless{}{-}} \FunctionTok{coef}\NormalTok{(mod.q4)[}\DecValTok{1}\NormalTok{] }\CommentTok{\# b0}
\NormalTok{  estimates\_c}\SpecialCharTok{$}\NormalTok{b1[r] }\OtherTok{\textless{}{-}} \FunctionTok{coef}\NormalTok{(mod.q4)[}\DecValTok{2}\NormalTok{] }\CommentTok{\# b1}
\NormalTok{\}}

\CommentTok{\# plot the estimates}
\NormalTok{p.b0.c }\OtherTok{\textless{}{-}} \FunctionTok{ggplot}\NormalTok{(estimates\_c, }\FunctionTok{aes}\NormalTok{(}\AttributeTok{x =}\NormalTok{ b0)) }\SpecialCharTok{+}
  \FunctionTok{geom\_histogram}\NormalTok{(}\FunctionTok{aes}\NormalTok{(}\AttributeTok{y =} \FunctionTok{after\_stat}\NormalTok{(density)), }\AttributeTok{bins =} \DecValTok{30}\NormalTok{,}
                 \AttributeTok{fill =} \StringTok{"lightgrey"}\NormalTok{, }\AttributeTok{color =} \StringTok{"black"}\NormalTok{) }\SpecialCharTok{+}
  \FunctionTok{geom\_density}\NormalTok{(}\AttributeTok{color =} \StringTok{"goldenrod"}\NormalTok{, }\AttributeTok{linewidth =} \DecValTok{1}\NormalTok{) }\SpecialCharTok{+}
  \FunctionTok{geom\_vline}\NormalTok{(}\AttributeTok{xintercept =} \DecValTok{2}\NormalTok{, }\AttributeTok{color =} \StringTok{"skyblue3"}\NormalTok{,}
             \AttributeTok{linetype =} \StringTok{"dashed"}\NormalTok{, }\AttributeTok{linewidth =} \DecValTok{1}\NormalTok{) }\SpecialCharTok{+}       \CommentTok{\# true beta\_0}
  \FunctionTok{labs}\NormalTok{(}\AttributeTok{x =} \FunctionTok{expression}\NormalTok{(}\FunctionTok{hat}\NormalTok{(beta)[}\DecValTok{0}\NormalTok{]), }\AttributeTok{y =} \StringTok{"Density"}\NormalTok{) }\SpecialCharTok{+}
  \FunctionTok{theme\_bw}\NormalTok{()}

\NormalTok{p.b1.c }\OtherTok{\textless{}{-}} \FunctionTok{ggplot}\NormalTok{(estimates\_c, }\FunctionTok{aes}\NormalTok{(}\AttributeTok{x =}\NormalTok{ b1)) }\SpecialCharTok{+}
  \FunctionTok{geom\_histogram}\NormalTok{(}\FunctionTok{aes}\NormalTok{(}\AttributeTok{y =} \FunctionTok{after\_stat}\NormalTok{(density)), }\AttributeTok{bins =} \DecValTok{30}\NormalTok{,}
                 \AttributeTok{fill =} \StringTok{"lightgrey"}\NormalTok{, }\AttributeTok{color =} \StringTok{"black"}\NormalTok{) }\SpecialCharTok{+}
  \FunctionTok{geom\_density}\NormalTok{(}\AttributeTok{color =} \StringTok{"goldenrod"}\NormalTok{, }\AttributeTok{linewidth =} \DecValTok{1}\NormalTok{) }\SpecialCharTok{+}
  \FunctionTok{geom\_vline}\NormalTok{(}\AttributeTok{xintercept =} \FloatTok{1.5}\NormalTok{, }\AttributeTok{color =} \StringTok{"skyblue3"}\NormalTok{,}
             \AttributeTok{linetype =} \StringTok{"dashed"}\NormalTok{, }\AttributeTok{linewidth =} \DecValTok{1}\NormalTok{) }\SpecialCharTok{+}       \CommentTok{\# true beta\_1}
  \FunctionTok{labs}\NormalTok{(}\AttributeTok{x =} \FunctionTok{expression}\NormalTok{(}\FunctionTok{hat}\NormalTok{(beta)[}\DecValTok{1}\NormalTok{]), }\AttributeTok{y =} \StringTok{"Density"}\NormalTok{) }\SpecialCharTok{+}
  \FunctionTok{theme\_bw}\NormalTok{()}

\NormalTok{p.b0.c }\SpecialCharTok{|}\NormalTok{ p.b1.c}
\end{Highlighting}
\end{Shaded}

\pandocbounded{\includegraphics[keepaspectratio]{HW3_files/figure-pdf/unnamed-chunk-25-1.pdf}}

\normalsize The values for \hat{\beta}\_0 and \hat{\beta}\_1 not
normally distributed around their true values of 2 and 1.5,
respectively. While most of the estimates seem to center around their
true values, it seems like the estimates that are different vary by a
lot more than the 2 scenarios above.

\section*{References}\label{bibliography}
\addcontentsline{toc}{section}{References}

\protect\phantomsection\label{refs}
\begin{CSLReferences}{1}{1}
\bibitem[\citeproctext]{ref-Bracht16}
{Bracht, Thilo, Christian Mölleken, Maike Ahrens, et al.} 2016.
{``Evaluation of the Biomarker Candidate {MFAP4} for Non-Invasive
Assessment of Hepatic Fibrosis in Hepatitis {C} Patients.''}
\emph{Journal of Translational Medicine} 14 (1): 201.

\bibitem[\citeproctext]{ref-Cui25}
Cui, Mengyang, and Jian Mao. 2025. {``Effect of Low Temperature on
Mechanical Properties and Fracture Behavior of 3D Printed Short Carbon
Fiber Reinforced Polyether Ether Ketone Composites (CF/PEEK).''}
\emph{Scientific Reports} 15 (1): 36244.

\bibitem[\citeproctext]{ref-interactions}
Long, Jacob A. 2024. \emph{Interactions: Comprehensive, User-Friendly
Toolkit for Probing Interactions}.
\url{https://doi.org/10.32614/CRAN.package.interactions}.

\end{CSLReferences}




\end{document}
